\documentclass[letterpaper, 12pt]{article}

\usepackage[utf8]{inputenc}
\usepackage[spanish]{babel} 
\usepackage{geometry} 

\geometry{margin=2.54cm}

\title{Introducción a la Contabilidad}
\author{Rosendo Camal}
\date{15 de marzo de 2026}

\pagestyle{plain}

\begin{document}

	\begin{titlepage}
		\centering
		{\scshape\LARGE Universidad Aztlán \par}
		\vspace{1cm}
		{\scshape\Large Sede Playa del Carmen \par}
		\vspace{3cm}
		{\scshape\Huge Fundamentos de Contabilidad \par}
		\vspace{1cm}
		{\itshape\Large Resumen de la Presentación \par}
		\vspace{3cm}
		{\Large José Rosendo Camal Teh \par}
		\vfill
		{\Large Ingeniería en Sistemas Computacionales \par}
		{\Large Segundo Cuatrimestre \par}
		{\Large Contabilidad Financiera \par}
		\vfill
		{\Large Profa. Guadalupe Parada Valencia \par}
		\vfill
		{\Large 15 de marzo de 2026 \par}
		\vfill
		{\Large Playa del Carmen, Quintana Roo \par}
	\end{titlepage}

	\section*{Introducción}
	
	Este documento resume la presentación sobre los \textbf{Fundamentos de Contabilidad} enfocado en el concepto, los tipos de contabilidad, los aspectos legales y los principios de contabilidad. El propósito es presentar la información de los temas mencionados de manera breve.
	
	\section*{Resumen de la presentación}

		\subsection*{Concepto de Contabilidad}
	
			La \emph{contabilidad} es una disciplina que se encarga de registrar, clasificar, guardar y elaborar reportes financieros de una entidad (persona, empresa, gobierno, etcétera) con el fin de permitir a la entidad o a sus interesados la mejor toma de decisiones desde una perspectiva financiera. Su objetivo general es registrar y clasificar transacciones, informar a través de documentos, cumplir con las normas legales y contables vigentes, dar el sustento y el apoyo a la toma de decisiones.

			Esta disciplina está sujeta a seguir unos principios basados en reconocer una entidad económica, en considerar que una transacción implica que la salida de una cuenta es la entrada de otra cuenta y viceversa, mantener el uso de una moneda base (de curso legal) en todos los registros de manera independiente de las divisas involucradas en las operaciones, especificar los periodos de la elaboración y presentación de la información contable y la importancia relativa de cada una de las transacciones llevadas a cabo. Este campo del saber se clasifica principalmente en los ámbitos financiero, administrativo, de costos y fiscal.
	
	\subsection*{Tipos de Contabilidad}

			La contabilidad se divide en distintas ramas, cada una con sus particulares pero todas centradas en registrar transacciones y presentar informes para la toma de decisiones para los interesados en dicha información y de una entidad en particular.

			\begin{itemize}

				\item \textbf{Contabilidad financiera:} Consiste en preparar la información contable de una entidad para sus interesados externos. La manera más común de visualizar el producto de este tipo de contabilidad es a través de los \emph{informes o kits para inversionistas}.
				\item \textbf{Contabilidad administrativa:} Es la rama encargada de reunir la información contable para ser usada por los directivos y gerentes pertenecientes a la entidad cuya contabilidad se elabora.
				\item \textbf{Contabilidad de costos:} Permite conocer los costos involucrados en la actividad de la entidad, establecer precios de los productos o servicios ofrecidos por dicha entidad, controlar los gastos y maximizar la utilidad.
				\item \textbf{Contabilidad fiscal:} Su objetivo es cumplir con la normativa legal y fiscal vigente, es decir, cumplir con las obligaciones fiscales de la entidad. La \emph{declaración de impuestos} es un trámite gubernamental que está dentro del alcance contable y es el ejemplo más conocido fuera de la contaduría fiscal.
				\item \textbf{Contabilidad gubernamental:} Es el tipo de contabilidad que está sujeto a mantener la información contable de las entidades de gobierno priorizando la transparencia y la rendición de cuenta. En la mayoría de los gobiernos de cualquier nivel (local o nacional) es visible el fruto de esta disciplina en sus informes de gobierno y en los presupuestos de proyectos y programas sociales, aunque esto es una versión limitada pero visible.
				\item \textbf{Contabilidad internacional:} Rama encargada de orquestar varias normas y prácticas legales de carácter local de una entidad que opera, desea o dejó de operar en dos o más países, con la finalidad de que la información contable sea uniforme y homogénea para ser usada en el contexto internacional (independiente de las prácticas y normas particulares locales).
				\item \textbf{Auditoría:} Consiste en supervisar que las prácticas y normas contables vigentes estén siendo llevadas a cabo de manera adecuada. Ayuda a los auditores a señalar los puntos de mejora en la contabilidad llevada por las entidades o a penalizar las prácticas siguiendo las mismas normas y prácticas que buscan proteger y según el objetivo de la auditoría.
			
			\end{itemize}

	\subsection*{Aspectos legales}
	
		Son los requisitos de la contabilidad que permiten, no solo la homogeneidad de la manera en la que se lleva a cabo la disciplina sino también la garantía de las auditorías, la presentación de la información contable, la transparencia y la facilidad en el cumplimiento de las obligaciones económicas siguiendo las pautas siguientes:
		
		\begin{itemize}

		\item Toda entidad está obligada y sujeta a llevar la contabilidad.
		\item Seguir la normativa vigente (\emph{NIIF} u otros según el caso).
		\item Cumplir con las obligaciones fiscales.
		\item La contabilidad se presenta en plazos trimestrales, semestrales y anuales (según el caso).
		\item Se conserva la contabilidad hasta cinco años (según sea el caso).
		\item Rendición de cuentas y regulación de auditorías.
		\item El personal a cargo, el contador, tiene responsabilidad legal en esta disciplina siendo sujeto a consecuencias marcadas por la ley.
		\item Cumplir con la ley laboral y la seguridad social.

		\end{itemize}
	
	\subsection*{Principios de Contabilidad}
	
		Son las reglas fundamentales que garantizan que las prácticas de la contabilidad sean consistentes, comprensibles y fiables. Pone en marcha un marco para la elaboración e interpretación de la información contable. Es decir, permite una estandarización de la disciplina. Esto permite un respaldo teórico del porqué se realiza la práctica de una manera y no otra con una justificación. Los principales principios son:
		
		\begin{itemize}

		\item \textbf{Entidad:} Es el que está sujeto a llevar la contabilidad, puede ser una persona o una empresa. En todo caso, la contabilidad de una entidad no debe mezclarse con otra entidad. El caso de error más común es juntar los gastos del dueño con su negocio.
		\item \textbf{Partida doble:} Principio que se resume como \emph{toda transacción afecta simultáneamente a dos cuentas}. La transacción de salida de una cuenta es la entrada de otra cuenta y viceversa.
		\item \textbf{Devengado:} Es la práctica de registrar cualquier movimiento en el momento en que se genera. Un caso particular a nivel persona física es comprar a crédito. Se debe anotar en nuestra contabilidad personal el gasto a pesar de que en el futuro se realice el pago. Esto permite tomar decisiones basadas en el presente inmediato evitando errores por considerar más dinero o menos del que se dispone y reserva el dinero para obligaciones no pagadas pero ya adquiridas.
		\item \textbf{Negocio en marcha:} Se considera que la empresa va a seguir funcionando en el futuro por lo que sus transacciones deben estar registradas a su valoración actual y no a precio de liquidación.
		\item \textbf{Periodicidad:} Periodo en el que se elabora y se presenta la información contable. Puede ser trimestral, semestral o anual según se requiera.
		\item \textbf{Moneda común:} Establecimiento de la moneda de curso legal que la empresa decida o deba usar en su contabilidad como unidad de medida, independientemente de las diferentes divisas involucradas en las operaciones de la entidad.
		\item \textbf{Costo histórico:} Los registros deben estar valuados al valor al que se efectuó, no al valor actualizado.
		\item \textbf{Relevancia:} Cada una de las transacciones tiene un peso en la contabilidad y la determinación de la importancia de cada una es vital para llevar a cabo la contabilidad. Por ejemplo, el costo de la adquisición de una oficina no tiene la misma relevancia que el gasto en papelería individual. A pesar de que ambos requieren ser registrados.
		\item \textbf{Consistencia:} Todas las transacciones, formatos, conceptos y plazos deben ser consistentes. Es decir, si un movimiento se registra como gasto corriente no puede volverse a registrar como costo sin una justificación. O bien, si se solicita un periodo trimestral, no puede ser entregado un informe contable de periodicidad anual.
		\item \textbf{Prudencia:} Se asegura primero los ingresos antes de anotarlos y se registran los gastos lo antes posible. Esto evita la generación errónea o pérdida de la información.
		\item \textbf{Objetividad:} Todo registro o transacción debe tener un respaldo comprobable y verificable. No debe haber juicios subjetivos. En otras palabras, cada gasto debe tener su comprobante de gasto y toda adquisición su comprobante de compra o contrato de adquisición, según sea el caso.
		\item \textbf{Universalidad:} Todos los principios se aplican a todas las entidades, sin embargo, cada una puede estar sujeta a pequeñas particularidades señaladas por la misma normativa vigente.

		\end{itemize}
	
		Estos principios permiten mejorar la transparencia de la información contable, poder comparar la información de la entidad con otras y de otros periodos de forma fiable y consistente y en general y como fin último, en la mejora de la toma de decisiones.
	
	\section*{Conclusión}
	
		En síntesis, la contabilidad es la disciplina cuya razón fundamental es mejorar la toma de decisiones mediante el uso de la información contable. Consiste en registrar las transacciones siguiendo una normatividad local o internacional, las prácticas de este campo del saber son estandarizadas para elevar la transparencia y rendición de cuentas ligadas al cumplimiento de las obligaciones económicas, laborales, fiscales y legales. Este trabajo permite conocer la contabilidad en sus términos generales lo que da espacio no solo para entender su finalidad sino en tener las bases para los temas siguientes. Sea el motivo de llevarlo a la práctica o al menos entender con mayor profundidad dicho campo del conocimiento.

\end{document}
